%───────────────────────────────────────────────────────────────────────────────
% METADATOS DEL DOCUMENTO — CURSO DATA STRUCTURES (CIIC3011 / ICOM4035)
%───────────────────────────────────────────────────────────────────────────────

% --- Identificación del documento ---
\newcommand{\docTitulo}{Estructuras de Datos}
\newcommand{\docSubtitulo}{De las estructuras lineales a las jerárquicas:\\fundamentos, algoritmos y complejidad}
\newcommand{\docAutor}{Jesús Antonio Chala Casanova}
\newcommand{\docFecha}{2026}
\newcommand{\docEstado}{Notas de Curso (CISE-Ph.D.)}

% --- Cita de portada ---
\newcommand{\docCita}{``La elección de una estructura de datos no es solo una decisión técnica; es una declaración sobre qué operaciones importan.''}

% --- Páginas preliminares ---
\newcommand{\docDedicatoria}{}
\newcommand{\docEpigrafe}{``Un algoritmo sin la estructura de datos correcta es como una ecuación diferencial sin condiciones de frontera: formalmente válido, prácticamente inútil.''}
\newcommand{\docAgradecimientos}{}

% --- Metadatos PDF ---
\newcommand{\docKeywords}{estructuras de datos, hashing, tablas hash, HashMap, TreeMap, árboles binarios, BST, árboles de búsqueda, recorridos, complejidad algorítmica, Java}
\newcommand{\docSubject}{Notas de estudio avanzado sobre estructuras de datos y algoritmos en Java}

% =============================================================================
% RECURSOS BIBLIOGRÁFICOS
% =============================================================================
\addbibresource{references/references.bib}
